\documentclass[11pt]{article}
\usepackage[margin=1in]{geometry}
\usepackage[T1]{fontenc}
\usepackage{lmodern}
\usepackage{amsmath,amssymb,amsthm}
\usepackage{microtype}
\usepackage[hidelinks]{hyperref}
\newtheorem{theorem}{Theorem}
\newtheorem{lemma}[theorem]{Lemma}
\newcommand{\R}{\mathbb R}
\newcommand{\C}{\mathbb C}
\newcommand{\nn}[1]{\left\lVert #1\right\rVert_2}
\title{Disproof of conjecture 00000001478\\
\large Polynomial images of cones and the tensor nuclear-norm ball}
\author{C0ldSmi1e}
\date{October 2026}
\begin{document}
\maketitle

\begin{abstract}
The nuclear-norm unit ball of complex \(2\times2\times2\) tensors cannot
be the image of a genuine cone under a polynomial map. A polynomial
bounded on a cone is constant along each ray, and its value on every
ray is its value at the common origin. The nuclear ball contains two
points distinguished by a bounded coordinate, giving a contradiction.
This rules out the conjectured spectrahedral-cone representation,
including the stated \(4\times4\) symmetric-matrix case.
\end{abstract}

\section{The source claim and the precise negation}

Both language versions of conjecture 00000001478 assert that the
nuclear-norm unit ball in \(\C^{2\times2\times2}\) is the polynomial
image of an explicit spectrahedral cone described using \(4\times4\)
symmetric matrices, and that this representation gives SDP computability
of the nuclear norm. The exact bilingual statement is included in
\texttt{ORIGINAL.md}.

We disprove the necessary representation assertion. Thus no interpretation
of the additional SDP consequence can make the stated conjunction true.
We do not claim that the nuclear norm cannot be computed by some other SDP
formulation.

The source abbreviates the nuclear norm as a \emph{minimal sum of rank-one
tensors}. We make its conventional scalar objective explicit. Let
\[
 V=\C^2,\qquad \mathcal T=\C^{2\times2\times2},
\]
with the usual Hermitian Euclidean norm on each factor \(V\).
For an arbitrary finite list \(L=((a_r,b_r,c_r))_r\) of triples in \(V^3\), put
\[
 A(L)_{ijk}=\sum_r (a_r)_i(b_r)_j(c_r)_k,\qquad
 c(L)=\sum_r \nn{a_r}\nn{b_r}\nn{c_r}.
\]
Define the standard complex projective tensor norm and its closed unit ball by
\[
 D(T)=\{c(L):A(L)=T\},\qquad
 \nu(T)=\inf D(T),\qquad B=\{T:\nu(T)\leq1\}.
\]
The lists in \(D(T)\) have unrestricted finite length. Every \(D(T)\) is
nonempty: the coordinate expansion
\[
 T=\sum_{i,j,k\in\{0,1\}}(T_{ijk}e_i)\otimes e_j\otimes e_k
\]
gives a decomposition. All costs are nonnegative.
Zero pure tensors can be deleted without changing either the sum or the
cost: a pure tensor is zero exactly when a factor is zero.
The formalization proves these facts, so allowing zero factors does not
change the attainable cost set for nonzero rank-one summands.

A cone here means a subset \(C\subseteq\R^n\) containing \(0\) and satisfying
\(tx\in C\) whenever \(x\in C\) and \(t\geq0\).
A polynomial map \(P:\R^n\to\mathcal T\) has real polynomials as the real
and imaginary parts of all eight complex coordinates.
Complex-coefficient polynomials on real parameters are included: their
real and imaginary parts are real polynomials, as also proved in Lean.

\begin{theorem}\label{thm:main}
For every natural number \(n\), every such cone \(C\subseteq\R^n\), and
every polynomial map \(P:\R^n\to\mathcal T\), one has \(P(C)\neq B\).
\end{theorem}

The quantifiers allow all finite parameter dimensions and all polynomial
degrees. Convexity and a particular matrix presentation are not needed.
This larger class contains every spectrahedral cone in the source claim.

\section{A bounded coordinate and two exact points}

\begin{lemma}\label{lem:bound}
For every tensor \(T\) and every coordinate \(i,j,k\),
\(\lvert T_{ijk}\rvert\leq\nu(T)\).
Consequently, \(\lvert\operatorname{Re}T_{000}\rvert\leq1\) for \(T\in B\).
\end{lemma}
\begin{proof}
For every factor triple,
\[
 \lvert a_i b_j c_k\rvert
 \leq \nn a\,\nn b\,\nn c.
\]
Applying the triangle inequality to any finite decomposition of \(T\)
gives \(\lvert T_{ijk}\rvert\leq c(L)\).
This is a lower bound for every element of the nonempty set \(D(T)\);
hence it is at most its infimum. The real-part bound follows from
\(\lvert\operatorname{Re}z\rvert\leq\lvert z\rvert\).
\end{proof}

Let \(E=e_0\otimes e_0\otimes e_0\).
The empty decomposition and nonnegative costs give \(\nu(0)=0\).
The one-term decomposition of \(E\) has cost \(1\), while
\(E_{000}=1\); Lemma~\ref{lem:bound} therefore gives \(\nu(E)=1\).
In particular,
\[
 0,E\in B,\qquad \operatorname{Re}0_{000}=0,\qquad
 \operatorname{Re}E_{000}=1.
\]
These are exact equalities, not numerical estimates.

\section{The polynomial obstruction}

\begin{lemma}\label{lem:ray}
If a real polynomial \(p\) on \(\R^n\) is bounded in absolute value on a
cone \(C\), then \(p(x)=p(0)\) for every \(x\in C\).
\end{lemma}
\begin{proof}
Fix \(x\in C\). Substitution \(X_i=x_i t\) gives the univariate polynomial
\(q_x(t)=p(tx)\). Cone closure makes \(q_x\) bounded for every \(t\geq0\).
A nonconstant real polynomial is unbounded in absolute value as
\(t\to+\infty\), by its leading term. Thus \(q_x\) is constant, and
\[
 p(x)=q_x(1)=q_x(0)=p(0).
\]
No homogeneity of \(p\), degree restriction, or interior point of \(C\)
is assumed.
\end{proof}

\begin{proof}[Proof of Theorem~\ref{thm:main}]
Suppose \(P(C)=B\). The real polynomial
\[
 p(x)=\operatorname{Re}(P(x)_{000})
\]
is bounded by \(1\) in absolute value on \(C\), by Lemma~\ref{lem:bound}.
Lemma~\ref{lem:ray} makes it equal to \(p(0)\) everywhere on \(C\).
The image equality supplies \(x_0,x_E\in C\) with \(P(x_0)=0\)
and \(P(x_E)=E\). Hence \(p(x_0)=0\) and \(p(x_E)=1\), a contradiction.
\end{proof}

Only one coordinate is needed for the contradiction.
More generally, coordinatewise application of Lemma~\ref{lem:ray}
shows that a bounded polynomial image of a nonempty cone is a singleton.

\section{The matrix condition and the word minimum}

For a real linear pencil \(L:\R^n\to\operatorname{Mat}_4(\R)\), define
\[
 C_L=\{x:L(x)\text{ is positive semidefinite}\}.
\]
The PSD condition includes symmetry. The usual pencils with symmetric
coefficient matrices are included in this definition.
One has \(0\in C_L\), since \(L(0)=0\). If \(x\in C_L\) and \(t\geq0\),
then \(L(tx)=tL(x)\) is symmetric and
\[
 v^{\mathsf T}L(tx)v=t\,v^{\mathsf T}L(x)v\geq0
 \quad\text{for all }v\in\R^4.
\]
Thus \(C_L\) is a cone, and Theorem~\ref{thm:main} rules out \(P(C_L)=B\).
Lean proves this closure directly from the genuine matrix PSD predicate.
The general theorem also rules out any further generalized matrix
description whose domain is an actual cone.

For completeness, the formalization separately treats the literal
attained-minimum reading
\[
 B_{\min}=\{T:\text{some }r\in D(T)\text{ is least in }D(T)
                         \text{ and }r\leq1\}.
\]
An attained least cost equals \(\inf D(T)\), so \(B_{\min}\subseteq B\).
The costs \(0\) for the zero tensor and \(1\) for \(E\) are attained
least costs, giving \(0,E\in B_{\min}\).
The proof of Theorem~\ref{thm:main} works unchanged for every set
\(B'\) with \(\{0,E\}\subseteq B'\subseteq B\), and therefore for
\(B_{\min}\). No unproved global attainment theorem is used.

The argument concerns the cone and polynomial map stated in both
languages. An affine section, normalization, or rational parametrization
would require a different statement.

\section{Formal verification and scope}

The project uses Lean 4.19.0 and Mathlib v4.19.0, pinned by the supplied
toolchain and dependency manifest.
\texttt{NuclearNorm.lean} defines the cost set and proves the tensor bounds
and exact points. \texttt{PolynomialCone.lean} proves the ray substitution,
constancy, and complex-coefficient correspondence.
\texttt{Solution.lean} combines them and proves the \(4\times4\) PSD case.
The main declaration is
\[
 \texttt{TLMC1478.conjecture00000001478\_false}.
\]
It is the negation of the existential representation by any genuine cone.
The whole source conjunction implies this existential assertion, so its
negation disproves the original conjecture.

The author build and strict source replays pass.
\texttt{AuthorAudit.lean} checks all explicitly declared objects and
results. The submission's additional verification records document
the independent rebuild, the audit of all compiled mathematical
declarations and their dependencies, and checks of the verification
tools themselves. Reproduction instructions and precise results are
in \texttt{README.md} and \texttt{VERIFICATION.md}.
Proofs may depend only on the standard foundational axioms
\texttt{propext}, \texttt{Classical.choice}, and \texttt{Quot.sound}.
No admitted proof, custom axiom, or \texttt{native\_decide} is used.

All mathematical computations are covered by the Lean proofs.
No numerical optimization, external symbolic calculation, or auxiliary
numerical program supplies a premise. The disproof needs neither all
abstract norm axioms nor global attainment: the standard cost formula
and the separately proved minimum-ball result suffice.
Verification here is local; acceptance remains a maintainer decision.
\end{document}
